% GENERATED FILE. Edit canonical lessons, then run npm run generate:books.
% canonical-source-hash: fnv1a64:8e2bda2b712e5ad0
% canonical-lessons: HI-C380-diqqat, HI-C380-badha, HI-C380-hal, HI-C380-atakna, HI-C380-ubarna

\chapter{Getting Out of Trouble}
\label{ch:hi-getting-out-of-trouble}
\hlchaptermodality{380}

\begin{chapteropening}
\textbf{By the end of this chapter:} \emph{I can say what to do if something gets in the way.}

Complete the last lesson of chapter 380: I can say what to do if something gets in the way.
\end{chapteropening}

\section[\texorpdfstring{\hi{दिक़्क़त} (diqqat)}{दिक़्क़त (diqqat)}]{\hi{दिक़्क़त} (diqqat) — a difficulty, a hassle}

\label{lesson:HI-C380-diqqat}



\begin{quote}
Before the new one: say the Hindi for beware, then the Hindi for to slam.
\end{quote}

\subsection*{You'll want to know: \hi{दिक़्क़त}}
\textbf{\hi{दिक़्क़त}} — \emph{diqqat} — \textquotedblleft{}a difficulty, a hassle\textquotedblright{}.

\begin{quote}
\hi{अगर} \hi{कोई} \hi{दिक़्क़त} \hi{होगी} \hi{तो} \hi{मुझे} \hi{फ़ोन} \hi{करो।}
\end{quote}

\begin{quote}
\emph{agar koī diqqat hogī to mujhe fon karo.}
\end{quote}

\begin{quote}
If there's any difficulty, phone me.
\end{quote}

Feminine.

\begin{cousinweb}
\textbf{\hi{दिक़्क़त}} (\emph{diqqat}) is usually traced to Arabic \emph{diqqa}, \textquotedblleft{}fineness, exactness\textquotedblright{}, which reached Hindi through Persian. The sense of a difficulty is a Persian and Hindi one.
\end{cousinweb}

\begin{grammarlens}[title={What you've built}]
One more word for saying what happens if.
\end{grammarlens}

\subsection*{Your turn}
\begin{itemize}
\raggedright
  \item \emph{Say it:} \emph{diqqat}
  \item \emph{Say it:} \emph{diqqat}, once more
  \item \emph{Say it:} \emph{agar koī diqqat hogī to mujhe fon karo.} — If there's any difficulty, phone me.
  \item \emph{Recall:} say the Hindi for beware, then the Hindi for to slam, then say \emph{diqqat} again
\end{itemize}

\begin{tcolorbox}[breakable,colback=teal!4,colframe=teal!35!black,title={Before you move on}]
What does \hi{दिक़्क़त} mean? (\textquotedblleft{}A difficulty, a hassle\textquotedblright{}.) Say it once more.
\end{tcolorbox}
\section[\texorpdfstring{\hi{बाधा} (bādhā)}{बाधा (bādhā)}]{\hi{बाधा} (bādhā) — an obstacle}

\label{lesson:HI-C380-badha}



\begin{quote}
Before the new one: say the Hindi for to slam, then the Hindi for a hassle.
\end{quote}

\subsection*{You'll want to know: \hi{बाधा}}
\textbf{\hi{बाधा}} — \emph{bādhā} — \textquotedblleft{}an obstacle\textquotedblright{}.

\begin{quote}
\hi{अगर} \hi{रास्ते} \hi{में} \hi{कोई} \hi{बाधा} \hi{आएगी} \hi{तो} \hi{हम} \hi{दूसरा} \hi{रास्ता} \hi{लेंगे।}
\end{quote}

\begin{quote}
\emph{agar rāste meṁ koī bādhā āegī to ham dūsrā rāstā leṅge.}
\end{quote}

\begin{quote}
If there's an obstacle on the way, we'll take another road.
\end{quote}

Feminine.

\begin{cousinweb}
\textbf{\hi{बाधा}} (\emph{bādhā}) is a tatsama, borrowed from Sanskrit \emph{bādhā}, \textquotedblleft{}an obstacle, a hindrance\textquotedblright{} (Monier-Williams).
\end{cousinweb}

\begin{grammarlens}[title={What you've built}]
One more word for saying what happens if.
\end{grammarlens}

\subsection*{Your turn}
\begin{itemize}
\raggedright
  \item \emph{Say it:} \emph{bādhā}
  \item \emph{Say it:} \emph{bādhā}, once more
  \item \emph{Say it:} \emph{agar rāste meṁ koī bādhā āegī to ham dūsrā rāstā leṅge.} — If there's an obstacle on the way, we'll take another road.
  \item \emph{Recall:} say the Hindi for to slam, then the Hindi for a hassle, then say \emph{bādhā} again
\end{itemize}

\begin{tcolorbox}[breakable,colback=teal!4,colframe=teal!35!black,title={Before you move on}]
What does \hi{बाधा} mean? (\textquotedblleft{}An obstacle\textquotedblright{}.) Say it once more.
\end{tcolorbox}
\section[\texorpdfstring{\hi{हल} (hal)}{हल (hal)}]{\hi{हल} (hal) — a solution}

\label{lesson:HI-C380-hal}



\begin{quote}
Before the new one: say the Hindi for a hassle, then the Hindi for an obstacle.
\end{quote}

\subsection*{You'll want to know: \hi{हल}}
\textbf{\hi{हल}} — \emph{hal} — \textquotedblleft{}a solution\textquotedblright{}.

\begin{quote}
\hi{अगर} \hi{हम} \hi{मिलकर} \hi{सोचेंगे} \hi{तो} \hi{कोई} \hi{हल} \hi{निकलेगा।}
\end{quote}

\begin{quote}
\emph{agar ham milkar soceṅge to koī hal niklegā.}
\end{quote}

\begin{quote}
If we think it over together, a solution will come.
\end{quote}

Masculine.

\begin{cousinweb}
\textbf{\hi{हल}} (\emph{hal}) is usually traced to Arabic \emph{ḥall}, \textquotedblleft{}loosening, solving\textquotedblright{}, which reached Hindi through Persian.
\end{cousinweb}

\begin{grammarlens}[title={What you've built}]
One more word for saying what happens if.
\end{grammarlens}

\subsection*{Your turn}
\begin{itemize}
\raggedright
  \item \emph{Say it:} \emph{hal}
  \item \emph{Say it:} \emph{hal}, once more
  \item \emph{Say it:} \emph{agar ham milkar soceṅge to koī hal niklegā.} — If we think it over together, a solution will come.
  \item \emph{Recall:} say the Hindi for a hassle, then the Hindi for an obstacle, then say \emph{hal} again
\end{itemize}

\begin{tcolorbox}[breakable,colback=teal!4,colframe=teal!35!black,title={Before you move on}]
What does \hi{हल} mean? (\textquotedblleft{}A solution\textquotedblright{}.) Say it once more.
\end{tcolorbox}
\section[\texorpdfstring{\hi{अटकना} (aṭaknā)}{अटकना (aṭaknā)}]{\hi{अटकना} (aṭaknā) — to get stuck, to be held up}

\label{lesson:HI-C380-atakna}



\begin{quote}
Before the new one: say the Hindi for an obstacle, then the Hindi for a solution.
\end{quote}

\subsection*{You'll want to know: \hi{अटकना}}
\textbf{\hi{अटकना}} — \emph{aṭaknā} — \textquotedblleft{}to get stuck, to be held up\textquotedblright{}.

\begin{quote}
\hi{अगर} \hi{दरवाज़ा} \hi{अटकेगा} \hi{तो} \hi{उसे} \hi{ज़ोर} \hi{से} \hi{धकेलो।}
\end{quote}

\begin{quote}
\emph{agar darvāzā aṭkegā to use zor se dhakelo.}
\end{quote}

\begin{quote}
If the door sticks, push it hard.
\end{quote}

Past (it …): \hi{वह} \hi{दरवाज़ा} \hi{अटक} \hi{गया}, \emph{vah darvāzā aṭak gayā}.

\begin{cousinweb}
\textbf{\hi{अटकना}} (\emph{aṭaknā}) is inherited from \emph{*aṭṭakk}, \textquotedblleft{}obstruct, stop\textquotedblright{}, a form CDIAL reconstructs from its descendants (CDIAL 182).
\end{cousinweb}

\begin{grammarlens}[title={What you've built}]
One more word for saying what happens if.
\end{grammarlens}

\subsection*{Your turn}
\begin{itemize}
\raggedright
  \item \emph{Say it:} \emph{aṭaknā}
  \item \emph{Say it:} \emph{aṭaknā}, once more
  \item \emph{Say it:} \emph{agar darvāzā aṭkegā to use zor se dhakelo.} — If the door sticks, push it hard.
  \item \emph{Recall:} say the Hindi for an obstacle, then the Hindi for a solution, then say \emph{aṭaknā} again
\end{itemize}

\begin{tcolorbox}[breakable,colback=teal!4,colframe=teal!35!black,title={Before you move on}]
What does \hi{अटकना} mean? (\textquotedblleft{}To get stuck, to be held up\textquotedblright{}.) Say it once more.
\end{tcolorbox}
\section[\texorpdfstring{\hi{उबारना} (ubārnā)}{उबारना (ubārnā)}]{\hi{उबारना} (ubārnā) — to pull out of trouble, to rescue}

\label{lesson:HI-C380-ubarna}



\begin{quote}
Before the new one: say the Hindi for a solution, then the Hindi for to get stuck.
\end{quote}

\subsection*{You'll want to know: \hi{उबारना}}
\textbf{\hi{उबारना}} — \emph{ubārnā} — \textquotedblleft{}to pull out of trouble, to rescue\textquotedblright{}.

\begin{quote}
\hi{अगर} \hi{मीरा} \hi{नहीं} \hi{होती} \hi{तो} \hi{मुझे} \hi{कौन} \hi{उबारता}?
\end{quote}

\begin{quote}
\emph{agar Mīrā nahīṃ hotī to mujhe kaun ubārtā?}
\end{quote}

\begin{quote}
If Mira hadn't been there, who would have pulled me out?
\end{quote}

Past, with \hi{ने} (he or she …): \hi{उसने} \hi{मुझे} \hi{मुसीबत} \hi{से} \hi{उबारा}, \emph{usne mujhe musībat se ubārā}.

\begin{cousinweb}
\textbf{\hi{उबारना}} (\emph{ubārnā}) is inherited from \emph{*udvārayati}, \textquotedblleft{}opens\textquotedblright{}, a form CDIAL reconstructs from its descendants (CDIAL 2082). CDIAL glosses the Hindi verb \textquotedblleft{}to release\textquotedblright{}.
\end{cousinweb}

\begin{grammarlens}[title={What you've built}]
One more word for saying what happens if.
\end{grammarlens}

\subsection*{Your turn}
\begin{itemize}
\raggedright
  \item \emph{Say it:} \emph{ubārnā}
  \item \emph{Say it:} \emph{ubārnā}, once more
  \item \emph{Say it:} \emph{agar Mīrā nahīṃ hotī to mujhe kaun ubārtā?} — If Mira hadn't been there, who would have pulled me out?
  \item \emph{Recall:} say the Hindi for a solution, then the Hindi for to get stuck, then say \emph{ubārnā} again
\end{itemize}

\begin{tcolorbox}[breakable,colback=teal!4,colframe=teal!35!black,title={Before you move on}]
What does \hi{उबारना} mean? (\textquotedblleft{}To pull out of trouble, to rescue\textquotedblright{}.) Say it once more.
\end{tcolorbox}
